\section{Hypothesis 3: the loop}
\label{sec:loop}

All results in this section use Qwen3-30B-A3B-Base unless stated, ten
narrative seeds, 4,500 tokens per cell, $\lambda=0$, and Claude Opus~5 as
judge with $k=5$. Full tables, including Cliff's $\delta$, Mann--Whitney
$p$-values and paired-by-seed intervals, are in Appendix~\ref{app:loop}.

\subsection{Bare generation is dead; the loop brings it to life}

Left to continue a premise with nothing else, the base model converges within
a few hundred tokens on an attractor of its pretraining and stays there:
\emph{``Highly recommended. Highly recommended.\,\ldots''} for two thousand
tokens; a table of three-digit numbers; an English-exam answer key. Judged
surprise 0.28, connection 0.20, coherence 2.06 (50 windows). The full scaffold
brings the stream to 3.19 / 1.64 / 3.68 (Cliff's $\delta=+0.88$ vs bare on
surprise, $p\approx5\times10^{-15}$; every seed moves in the same direction,
Figure~\ref{fig:perseed}), and elaborates its own finds after a
re-encounter (\emph{``Everything that almost happened piled up inside her like
snow falling sideways through open windowsills into rooms filled with silence
instead of noise\,\ldots\ She kept a notebook because sometimes almosts piled
up too much to be ignored''}). Salience-gated review also beats reviewing on a
timer as a policy for \emph{which} windows to judge (surprise and connection
CIs positive at a third of the judge cost).

\begin{figure}[t]
\centering
\includegraphics[width=0.6\linewidth]{fig3_per_seed_scaffold_vs_bare.png}
\caption{Per-seed robustness: mean surprise of bare vs scaffold windows, one
line per seed.}
\label{fig:perseed}
\end{figure}

\subsection{Taking the loop apart}
\label{sec:ablation}

Ablation over 10 seeds (Figure~\ref{fig:ablation}, Table~\ref{tab:ladder})
gave the first surprise: \emph{bare + clock reseed} (no salience, no
forgetting, no kicks, only a neutral subject change every 150 tokens over the
preserved context) ties the full scaffold on surprise (3.08 vs 3.19, CI
\ci{-0.54}{+0.78}) and beats it on connection (3.15 vs 1.64, $\delta=-0.58$
for the scaffold, $p\approx1.5\times10^{-7}$) and coherence (4.78 vs 3.68), with
12/60 windows judged both surprising and coherent against 8/47 for the
scaffold and 0/50 for bare. Salience-only (review windows marked but nothing
injected) reaches 2.63. In sentence-embedding space
(Figure~\ref{fig:traj}) bare generation has explored radius 0.18 and mean step
0.14 and ends 0.78 from the premise; the scaffold has radius 0.51 and returns
to within 0.57 of the premise; the clock reseed spreads over the space in
jumps while keeping the whole context, which is where its connection
advantage comes from.

\begin{figure}[t]
\centering
\includegraphics[width=\linewidth]{fig1_reverie_distributions.png}
\caption{Ablation: judged surprise, connection and coherence by arm (10
seeds, Opus~5, median of $k=5$ per window; violins are distributions, bars are
means with bootstrap CIs). Bare + clock reseed matches the full scaffold on
surprise and exceeds it on connection and coherence.}
\label{fig:ablation}
\end{figure}

\begin{figure}[t]
\centering
\includegraphics[width=\linewidth]{traj_dream_scaffold_s6.png}
\caption{Where the stream goes: trajectories in sentence-embedding space
(64-token windows, one PCA per seed), seed 6. $\star$ premise, $\circ$ start,
$\times$ reseed, $\diamond$ judged event sized by surprise. Bare makes one
jump and freezes; the scaffold orbits and returns; the clock reseed jumps
across the space keeping its context.}
\label{fig:traj}
\end{figure}

\subsection{What the interruption is made of}
\label{sec:battery2}

The ablation left the effect in one operation and raised four questions about
it, and one about itself. Battery 2 answers them (Table~\ref{tab:ladder},
Figures~\ref{fig:q1}--\ref{fig:q4}).

\paragraph{A confound, and two operations.} The bare arm of the ablation ran
without the scaffold's habituation penalty, which every interrupted arm had;
bare generation dies in \emph{literal} orbits, precisely what a repetition
penalty undoes. Adding the missing control (habituation, no interruption)
splits the effect in two: bare 0.28 $\rightarrow$ bare + habituation
\textbf{1.70} (surprise $\Delta$ +1.42 \ci{+1.12}{+1.72} paired by seed;
connection 0.20 $\rightarrow$ 1.08; coherence 2.06 $\rightarrow$ 4.08)
$\rightarrow$ bare + habituation + clock reseed \textbf{3.08} ($\Delta$ +1.38
\ci{+0.92}{+1.92}). Habituation removes the literal orbits and buys about
half of the surprise; the interruption buys the other half and \emph{all} of
the connection (1.08 $\rightarrow$ 3.15, $\Delta$ +2.07 \ci{+1.63}{+2.57},
$\delta=+0.75$) and the coherence. The full scaffold adds surprise +1.49
\ci{+0.80}{+2.19} over habituation alone but only +0.55 of connection.

\begin{table}[t]
\centering\small
\begin{tabular}{@{}lccc@{\hspace{1.2em}}ccc@{\hspace{1.2em}}ccc@{}}
\toprule
& \multicolumn{3}{c}{Qwen3-30B-A3B} & \multicolumn{3}{c}{Qwen3-8B} & \multicolumn{3}{c}{OLMo-2-13B} \\
\cmidrule(lr){2-4}\cmidrule(lr){5-7}\cmidrule(lr){8-10}
arm & S & C & H & S & C & H & S & C & H \\
\midrule
bare & 0.28 & 0.20 & 2.06 & 0.68 & 0.38 & 3.45 & 1.35 & 0.91 & 2.97 \\
bare + habituation & 1.70 & 1.08 & 4.08 & 1.47 & 0.87 & 3.93 & 2.37 & 1.51 & 4.05 \\
bare + habituation + clock reseed & \textbf{3.08} & \textbf{3.15} & \textbf{4.78} & \textbf{2.73} & \textbf{3.23} & 4.42 & \textbf{3.11} & \textbf{3.17} & 3.28 \\
DREAM scaffold & 3.19 & 1.64 & 3.68 & 2.73 & 1.53 & 4.30 & 3.08 & 1.89 & \textbf{4.48} \\
\bottomrule
\end{tabular}
\caption{The ladder on three generator families: judged surprise (S),
connection (C) and coherence (H), means over windows (Opus~5, $k=5$; 10 seeds
per arm; windows before 100 generated tokens excluded). Habituation buys part
of the surprise; the interruption buys the rest and all of the connection; the
scaffold adds nothing over the plain interruption except, on OLMo, coherence
(its selective forgetting lifts the stream out of a web-boilerplate well).}
\label{tab:ladder}
\end{table}

\begin{figure}[t]
\centering
\includegraphics[width=\linewidth]{fig8_q1_habituation.png}
\caption{Habituation or interruption? bare $\rightarrow$ bare + habituation
$\rightarrow$ bare + habituation + clock reseed $\rightarrow$ scaffold
(Qwen3-30B-A3B).}
\label{fig:q1}
\end{figure}

\paragraph{The interruption must lead away.} With the same period (150), a
neutral subject change and a re-encounter stitch that asks the stream to
return to its beginning are close (3.08 / 3.15 / 4.78 vs 2.68 / 3.64 / 4.72;
the stitch trades some surprise for connection, not significantly). Injecting
the \emph{premise itself} (1.14 / 1.04 / 3.18) or a window of the stream's
\emph{own past} (1.42 / 1.52 / 3.06) is as bad as not interrupting at all:
every paired CI against the neutral change excludes zero (connection $-2.11$
\ci{-2.83}{-1.46} and $-1.63$ \ci{-2.37}{-0.91}; $\delta\approx-0.6$ to
$-0.7$). Qualitatively, the own-past injection reinforces whatever attractor
the stream is in: an exam-key well was fed its own questions back. A return
works only as a short stitch that still opens a new sentence; a literal return
closes the loop (Figure~\ref{fig:q2}).

\begin{figure}[t]
\centering
\includegraphics[width=\linewidth]{fig8_q2_content.png}
\caption{Content of the interruption at period 150: neutral change,
re-encounter stitch, the premise itself, the stream's own past. Injecting a
return is as bad as not interrupting.}
\label{fig:q2}
\end{figure}

\paragraph{Salience is a good reader and a bad metronome.} The stitch
injected on salience events (jump, crystallization, recurrence; 1--9 per cell,
median $\approx$5) makes the \emph{event windows} better than salience-only
(3.92 / 2.99 / 3.75 vs 2.63 / 1.39 / 3.71), but the \emph{stream} it produces
is poor: uniform windows 2.18 / 1.68 / 3.93 against \textbf{4.70 / 3.48 /
5.50} for the same stitch on a clock of matched frequency (every 900 tokens;
surprise $\Delta$ +2.02 \ci{+1.62}{+2.44} paired, coherence +0.78) and 5.48 /
2.87 / 5.95 for a neutral change every 900 (45/60 windows both surprising and
coherent, the best cell of the program). Salience events select good moments
to look at, which is why salience-gated review beat clock review; as an
interruption trigger they fire rarely and unevenly and leave streams stuck for
thousands of tokens. Salience should decide where to look, not when to
interrupt (Figure~\ref{fig:q3}).

\begin{figure}[t]
\centering
\includegraphics[width=\linewidth]{fig8_q3_timing.png}
\caption{When to return: the re-encounter stitch on the clock (150 and 900),
on salience events, and salience-only; event windows and uniform windows.}
\label{fig:q3}
\end{figure}

\paragraph{The yield of an interruption depends on the thread it breaks, and
decays with distance from it.} With the neutral change every 75 tokens the
stream is dead (0.88, $\delta=-0.83$ vs 150): nothing has time to develop.
Every 150: 3.08. Every 300, 600 and 900, the windows that end within
$\approx$160 tokens after an interruption score \textbf{5.28 / 3.35 / 5.95},
5.40 / 3.20 / 5.65 and 5.90 / 2.80 / 6.10, the highest values in the
program (34/60 of the 300-early windows both surprising and coherent), while
windows deeper into a segment fall back to $\approx$3.2 and stay there. The
same window shape (a break plus 150 tokens of development) scores 3.08 when
the broken thread was 150 tokens long and 5.3--5.9 when it was 300--900:
surprise needs a developed background to break. Weighting the two phases by
their share of the stream, period 300 gives the best stream (4.35 / 3.18 /
5.97), then 600 (3.82), 900 (3.68), 150 (3.08) and 75 (0.88); connection
falls as segments lengthen (Figure~\ref{fig:q4}). The loop's rhythm is: let a
thread develop for a few hundred tokens, then break it away.

\begin{figure}[t]
\centering
\includegraphics[width=\linewidth]{fig8_q4b_phase.png}
\caption{How often to interrupt. Left: judged surprise against tokens since
the last interruption, per period. Right: stream mean by period. Surprise
peaks in the 150 tokens after a break and decays; a break every 75 tokens
develops nothing.}
\label{fig:q4}
\end{figure}

\subsection{Three generator families}
\label{sec:families}

The ladder replicates on Qwen3-8B-Base and OLMo-2-13B (Table~\ref{tab:ladder},
Figures~\ref{fig:fam8b}--\ref{fig:famolmo}). On the 8B: 0.68 $\rightarrow$
1.47 $\rightarrow$ 2.73 (connection 0.87 $\rightarrow$ 3.23, $\delta=+0.80$),
scaffold 2.73 / 1.53. On OLMo, whose bare stream wanders across web genres
rather than locking into literal orbits and is therefore less dead
(1.35), the same ordering holds on surprise and connection (2.37
$\rightarrow$ 3.11; connection 1.51 $\rightarrow$ 3.17, $\delta=+0.58$), with
one new nuance: the plain interruption \emph{costs} coherence there ($-0.77$
\ci{-1.58}{+0.03}; the injected narrative seed lands on 4,000 tokens of
footers) and the scaffold's selective forgetting is what preserves it (4.48
vs 3.28), the mechanism the calibration probes had suggested. Whether
forgetting is worth its connection cost depends on how deep the generator's
well is.

\begin{figure}[t]
\centering
\includegraphics[width=\linewidth]{fig8_q5_family8b.png}
\caption{The ladder on Qwen3-8B-Base.}
\label{fig:fam8b}
\end{figure}
\begin{figure}[t]
\centering
\includegraphics[width=\linewidth]{fig8_q6_familyolmo.png}
\caption{The ladder on OLMo-2-13B: same ordering; the plain interruption
costs coherence and forgetting preserves it.}
\label{fig:famolmo}
\end{figure}

